\documentclass[a4paper,11pt]{article}

\usepackage[utf8]{inputenc}
\usepackage[T1]{fontenc}
\usepackage{geometry}
\usepackage{tabularx}
\usepackage{booktabs}
\usepackage{array}
\usepackage{xcolor}
\usepackage{amsmath}
\usepackage{ragged2e}
\usepackage[hidelinks]{hyperref}

\geometry{
    a4paper,
    top=2cm,
    bottom=2cm,
    left=2cm,
    right=2cm
}

\setlength{\parindent}{0pt}
\setlength{\parskip}{6pt}
\renewcommand{\arraystretch}{1.25}

\newcolumntype{L}[1]{>{\RaggedRight\arraybackslash}p{#1}}
\newcolumntype{Y}{>{\RaggedRight\arraybackslash}X}

\title{\textbf{DEOS Sistem Mimarisi: Fonksiyonel Dağılım ve Sorumluluk Raporu}\\
\large Sürüm 0.1 (Taslak) -- Otonom, Merkezi Beyin ve Alt ECU Sorumluluk Matrisi}
\author{\textbf{Emir Arkalı} (\href{mailto:emirarkali07@gmail.com}{emirarkali07@gmail.com})\\DEOS -- Dokuz Eylül Otonom Sistemler}
\date{İlk Yayın: 28 Eylül 2026 \quad | \quad Son Revizyon: 13 Ekim 2026}

\begin{document}

\maketitle

\begin{center}
\small
\begin{tabularx}{\textwidth}{p{3.2cm}p{4.5cm} p{3.2cm}X}
\toprule
\textbf{Doküman Tipi:} & Mimari Tasarım Raporu & \textbf{Doküman Kodu:} & \texttt{DOC-REP-SPLIT-001} \\
\textbf{Doküman Sürümü:} & Sürüm 0.1 (Taslak / Draft) & \textbf{Sistem Sürümü:} & DEOS v1.1 Mimarisi \\
\textbf{Yazar:} & Emir Arkalı & \textbf{İnceleyen:} & DEOS Mühendislik Ekibi \\
\textbf{İlk Yayın Tarihi:} & 28 Eylül 2026 & \textbf{Son Revizyon Tarihi:} & 13 Ekim 2026 \\
\textbf{Durum:} & İnceleniyor / Erken Taslak & \textbf{Kapsam:} & ROS 2, Bridge, Main, Alt ECU'lar \\
\bottomrule
\end{tabularx}
\end{center}

\vspace{1em}

\section{Giriş ve Problem Tanımı}

Karmaşık dağıtık gömülü sistemlerde ve otonom araç platformlarında, sistemin başarısı yalnızca haberleşme paketlerinin bayt seviyesinde tanımlanmasıyla (ICD) sınırlı değildir. En sık karşılaşılan sistem entegrasyonu problemleri, \textbf{"Hangi fonksiyonun nerede koştuğu"} konusundaki gri alanlardan (belirsizliklerden) kaynaklanır.

Örneğin:
\begin{itemize}
    \item Hız ve tekerlek açısı PID hesaplamaları otonom yazılımda mı yoksa motor sürücüde mi koşmalıdır?
    \item Sinyal lambasının yanıp sönme frekansı (1 Hz flaşör) otonom planlayıcı tarafından mı tetiklenmelidir, yoksa donanımsal timer'a mı devredilmelidir?
    \item Acil durdurma ve timeout durumlarında aracı durdurma yetkisi kimdedir?
\end{itemize}

Bu rapor; araç üstü **ROS 2 Otonom Sürüş Bilgisayarı (IPC)**, **Ağ Köprüsü (\texttt{deos\_ros\_bridge})**, **Merkezi Araç Kontrolcüsü (\texttt{Main STM32 / VCU})** ve **Alt Donanım ECU'ları (Traction, Steering, Brake, BMS)** arasındaki görev paylaşımını ve yetki hiyerarşisini (Functional Split) resmiyete döker.

\section{Mimari Katmanlar ve Düğüm Görevleri}

Sistem hiyerarşik olarak 4 ana operasyonel katmana ayrılmıştır:

\begin{enumerate}
    \item \textbf{Katman 1: Üst Seviye Otonom Sürüş (ROS 2 / IPC)}:
    Çevre algılama (LiDAR, Kamera), yapay zeka nesne tespiti, SLAM haritalama, rota ve yörünge planlaması. Saf SI birimleriyle ($m/s$, $rad$, $rad/s$) çalışır.
    
    \item \textbf{Katman 2: Ağ Köprüsü (\texttt{deos\_ros\_bridge} / UDP Port 5001)}:
    ROS 2 SI birimleri ile gömülü ağın fixed-point ($0.01\ m/s$, $0.01^\circ$) değerleri arasındaki köprüdür. Ackermann kinematiğini çözer ve 200 ms komut zaman aşımını takip eder.
    
    \item \textbf{Katman 3: Merkezi Araç Beyni (\texttt{Main STM32 / VCU})}:
    Aracın durum makinesini (State Machine) yönetir. Donanımsal E-STOP butonu, ağ yönlendiricisi (Router), ağ canlılık bekçisi (Watchdog) ve aydınlatma donanım sürücüsünü (\texttt{LedController}) barındırır.
    
    \item \textbf{Katman 4: Aktüatör ve Sensör Düğümleri (CAN-FD Omurgası)}:
    Traction ECU (Motor akım ve PWM sürücüsü), Steering ECU (Direksiyon kapalı çevrim kontrolü), Brake ECU (Hidrolik/mekanik frenleme) ve BMS Gateway (Batarya hücre koruması).
\end{enumerate}

\clearpage

\section{Fonksiyonel Sorumluluk Matrisi (Functional Allocation)}

Aşağıdaki tablo, sistemdeki temel işlevlerin hangi katmanlar tarafından yürütüldüğünü kesin çizgilerle tanımlar:

\begin{table}[h]
\centering
\small
\begin{tabularx}{\textwidth}{|L{4.2cm}|L{2.8cm}|L{2.8cm}|L{2.8cm}|Y|}
\hline
\textbf{Fonksiyonel Alan} & \textbf{ROS 2 Otonom} & \textbf{ROS 2 Bridge} & \textbf{Main STM32} & \textbf{Alt ECU'lar} \\
\hline
\textbf{Algılama \& SLAM} & \textbf{Tam Sorumlu} & İlgisiz & İlgisiz & İlgisiz \\
\hline
\textbf{Yörünge Planlama} & \textbf{Tam Sorumlu} (\texttt{cmd\_vel}) & İlgisiz & İlgisiz & İlgisiz \\
\hline
\textbf{Ackermann Kinematiği} & İlgisiz & \textbf{Çevirici} ($v, \omega \rightarrow v, \delta$) & Limit Denetimi & İlgisiz \\
\hline
\textbf{Birim Ölçekleme} & İlgisiz (Saf SI) & \textbf{Tam Sorumlu} & Fixed-point & Fixed-point \\
\hline
\textbf{Motor Akım/Hız PID} & İlgisiz & İlgisiz & İlgisiz & \textbf{Traction ECU} \\
\hline
\textbf{Direksiyon Açı PID} & İlgisiz & İlgisiz & İlgisiz & \textbf{Steering ECU} \\
\hline
\textbf{Fren Basıncı \& EBS} & İlgisiz & İlgisiz & E-STOP Talebi & \textbf{Brake ECU} \\
\hline
\textbf{Sinyal Flaşörü (1 Hz)} & Niyet Bildirimi & İletici & \textbf{Donanım Timer} & İlgisiz \\
\hline
\textbf{Fiziksel E-STOP} & İlgisiz & İlgisiz & \textbf{En Üst Otorite} & STO Koruması \\
\hline
\textbf{Araç Durum Makinesi} & Durum İzleyici & İletici & \textbf{Tek Karar Verici}& Durum Takipçisi \\
\hline
\textbf{Batarya Koruma / BMS} & Telemetri İzleyici & İletici & Güvenlik Kapatması & \textbf{BMS Gateway} \\
\hline
\end{tabularx}
\caption{DEOS Fonksiyonel Dağılım ve Görev Matrisi}
\end{table}

\section{Kritik Sınırların ve Gri Alanların Çözümü}

\subsection{Aydınlatma ve Sinyal Lambası Ayrımı (Blinker Logic)}
\begin{itemize}
    \item Otonom sürüş yazılımı saniyede 1 kez lambayı açıp kapatarak ağ trafiğini gereksiz paketlerle boğmaz.
    \item ROS 2 sadece niyet bildirir: \texttt{/vehicle/signals/right = true} (\texttt{On Change}).
    \item Main STM32 üzerindeki \texttt{LedController}, dahili 10 ms döngüsünde çalışan zamanlayıcı ile lambayı 1 Hz periyotla (\%50 duty cycle) donanımsal olarak yakıp söndürür. Sürüş tamamlandığında ROS 2 \texttt{false} gönderir ve donanım söner.
\end{itemize}

\subsection{Kinematik Dönüşüm ve Motor Kontrol Sınırı}
\begin{itemize}
    \item ROS 2 aracı kinematik bir model olarak görür ve hedef hız ($v = \text{linear.x}$) ile hedef yaw açısal hızı ($\omega = \text{angular.z}$) üretir.
    \item \texttt{deos\_ros\_bridge}, aracın dingil mesafesini ($L$) kullanarak Ackermann tekerlek açısını hesaplar:
    $$\delta = \arctan\left(\frac{L \cdot \omega}{v}\right)$$
    \item Elde edilen hedef hız Traction ECU'ya, hedef açı Steering ECU'ya aktarılır.
    \item Her aktüatör ECU'su kendi içindeki yüksek frekanslı (1 kHz) akım ve tork geri besleme döngüsünü bağımsız işletir.
\end{itemize}

\subsection{Yetki Hiyerarşisi (Command Arbitration)}
Hiçbir koşulda otonom planlayıcı doğrudan tekerleklere müdahale edemez. Çok katmanlı hakemlik kuralı uygulanır:

\vspace{0.3em}
\begin{center}
\textbf{E-STOP (Donanım)} $\boldsymbol{\gg}$ \textbf{Yer Kontrolü (LoRa)} $\boldsymbol{\gg}$ \textbf{Watchdog Hatası} $\boldsymbol{\gg}$ \textbf{Manuel} $\boldsymbol{\gg}$ \textbf{Otonom (ROS 2)}
\end{center}
\vspace{0.3em}

\begin{itemize}
    \item Fiziksel E-STOP butonuna basıldığında Main STM32 \texttt{SafetyManager} anında \texttt{DEOS\_STATE\_SAFE} moduna geçer.
    \item Bu komut CAN-FD omurgasında en yüksek öncelikle (\texttt{DEOS\_PRIO\_EMERGENCY}) yayınlanır.
    \item ROS 2'nin bu güvenli duruş durumunu ezme veya devre dışı bırakma yetkisi kesinlikle yoktur.
\end{itemize}

\clearpage

\section{Zaman Aşımı ve Failsafe Matrisi}

Sistem bileşenlerinin iletişim kaybı veya donanım arızası yaşaması durumunda devreye girecek deterministik koruma mekanizmaları:

\begin{table}[h]
\centering
\small
\begin{tabularx}{\textwidth}{|L{4.5cm}|L{3.0cm}|c|Y|}
\hline
\textbf{Hata / Arıza Durumu} & \textbf{Tespit Eden Katman} & \textbf{Süre} & \textbf{Failsafe Davranışı} \\
\hline
\textbf{ROS 2 Komut Kesilmesi} & \texttt{deos\_ros\_bridge} & 200 ms & Hedef hız sıfırlanır, yumuşak duruş talep edilir. \\
\hline
\textbf{Ethernet Soket Kopması} & Main STM32 (UDP/TCP) & 500 ms & Otonom izin iptal edilir, araç \texttt{STANDBY}'a geçer. \\
\hline
\textbf{Traction ECU İletişim Kaybı} & Main STM32 (CAN-FD) & 500 ms & \texttt{DEOS\_STATE\_SAFE} ilan edilir, acil fren kilitlenir. \\
\hline
\textbf{Steering Açı Sensör Hatası} & Steering ECU & $<50$ ms & Motor torku kesilir, hata kodu fırlatılır. \\
\hline
\textbf{BMS Aşırı Sıcaklık/Voltaj} & BMS Gateway & $<100$ ms & Kritik arıza bildirilir, kontaktör kontrollü açılır. \\
\hline
\textbf{Fiziksel E-STOP Butonu} & Main STM32 GPIO & $<1$ ms & Tüm çekiş motorları enerjisiz kalır (STO), EBS kilitlenir. \\
\hline
\end{tabularx}
\caption{Failsafe ve Zaman Aşımı Tetikleyicileri}
\end{table}

\section{Sonuç ve Geliştirme Yol Haritası}

Bu rapor, sistemin erken geliştirme aşaması için bir taslak (\texttt{v0.1-DRAFT}) teşkil etmektedir. Projenin sonraki fazlarında:
\begin{enumerate}
    \item Dinamik tekerlek kayma açısı kompanzasyonu (Slip angle compensation),
    \item Çift motorlu araçlar için Tork Vektörleme (Torque Vectoring) mantığının Traction ECU'ya yerleştirilmesi,
    \item Mesafe sensörleri ile acil otonom frenleme (AEB) eşikleri,
\end{enumerate}
ayrıntılı olarak bu rapora dahil edilecektir.

\end{document}
